\chapter{锅炉设备}

\section{思考题}

\begin{problem}
何谓锅炉设备？其各组成设备的作用是什么？简述锅炉设备的工作过程。
\end{problem}

\begin{solution}
\begin{enumerate}
  \item \textbf{锅炉设备的定义}：锅炉设备是将燃料的化学能转变为工质热能的设备，它是锅炉本体设备及其辅助设备的总称。
  \item \textbf{各组成设备的作用}：
  \begin{enumerate}
    \item \textbf{"锅"（汽水系统）}：用于盛放工质，并接受燃料燃烧放出的热量，使工质由水逐渐被加热至合乎要求的过热蒸汽。
    \item \textbf{"炉"（燃烧系统）}：使燃料在其内部快速、稳定、完全地燃烧，放出热量，产生高温火焰和烟气。
    \item \textbf{辅助设备}：保障锅炉本体的正常运转，主要包括给水设备、通风设备、输煤设备、制粉设备、除尘及除灰设备等。
  \end{enumerate}
  \item \textbf{锅炉设备的工作过程}：
  \begin{enumerate}
    \item \textbf{燃烧过程}：制粉系统磨制的煤粉经一次风送入炉膛，同时送入二次风助燃。气粉混合物在炉膛内悬浮燃烧放出热量，高温火焰及烟气向各受热面进行辐射和对流换热，最后烟气经除尘脱硫由引风机排入大气。
    \item \textbf{产汽过程}：给水经省煤器加热后进入汽包，再经下降管到水冷壁吸收辐射热而部分汽化；汽水混合物回到汽包后分离出饱和蒸汽，饱和蒸汽再经过热器加热为过热蒸汽送往汽轮机做功。
  \end{enumerate}
\end{enumerate}
\end{solution}

\begin{problem}
煤的元素分析成分中哪些是可燃成分？哪些是不可燃成分？其中硫和氮的氧化物对设备和环境的影响如何？
\end{problem}

\begin{solution}
\begin{enumerate}
  \item \textbf{可燃成分}：碳(C)、氢(H)、硫(S)。
  \item \textbf{不可燃成分}：氧(O)、氮(N)、灰分(A)、水分(M)。
  \item \textbf{硫和氮氧化物的影响}：
  \begin{enumerate}
    \item \textbf{硫(S)}：燃烧产物为 \(\text{SO}_2\) 或者 \(\text{SO}_3\)，会污染大气，并使设备发生腐蚀。
    \item \textbf{氮(N)}：为有害元素，其氧化产物 \(\text{NO}\) 或 \(\text{NO}_2\) 的排放会对环境有害。
  \end{enumerate}
\end{enumerate}
\end{solution}

\begin{problem}
何谓挥发分、发热量及标准煤？电厂燃煤一般分为几种？其挥发分和发热量如何？
\end{problem}

\begin{solution}
\begin{enumerate}
  \item \textbf{定义}：
  \begin{enumerate}
    \item \textbf{挥发分}：煤在加热过程中有机质分解而析出的气体物质。
    \item \textbf{发热量}：单位质量收到基燃料完全燃烧时释放的热量。
    \item \textbf{标准煤}：规定低位发热量 \(Q_{ar,net,p}=\qty{29308}{kJ/kg}\) 的煤作为标准煤。
  \end{enumerate}
  \item \textbf{电厂燃煤分类及特点}：
  \begin{enumerate}
    \item \textbf{无烟煤}：挥发分含量最低（\(<10\%\)），发热量较高。
    \item \textbf{贫煤}：挥发分较低（\(10\%\sim20\%\)），发热量较高。
    \item \textbf{烟煤}：挥发分较高（\(20\%\sim40\%\)），发热量也相当高。
    \item \textbf{褐煤}：挥发物含量高（\(>40\%\)），易于着火，但灰分高、发热量较低。
    \item （此外还包括煤泥、泥煤等低质煤，其发热量很低，灰分极高）。
  \end{enumerate}
\end{enumerate}
\end{solution}

\begin{problem}
煤粉有哪些特性？制粉系统分为哪些类型？不同类型的制粉系统一般配备哪种磨馍机？不同型式的磨煤机的特点如何？
\end{problem}

\begin{solution}
\begin{enumerate}
  \item \textbf{煤粉的特性}：具有良好的流动性、易发生自燃及爆炸，并对燃烧有直接影响的煤粉细度（越细越易完全燃烧，但制粉电耗增加）。
  \item \textbf{制粉系统类型}：分为直吹式制粉系统和中间储仓式制粉系统。
  \item \textbf{配备磨煤机及特点}：
  \begin{enumerate}
    \item \textbf{直吹式制粉系统}：一般配用\textbf{中速磨煤机或风扇式（高速）磨煤机}。特点：系统结构简单、布置紧凑、占地少、初投资小和运行电耗低，但对磨煤机的可靠性要求较高。
    \item \textbf{中间储仓式制粉系统}：一般配用\textbf{低速滚筒式钢球磨煤机}。特点：钢球磨对煤种适应范围广，可以磨制任何难磨的煤种；但系统结构庞大复杂、占地面积大、初投资大，且电耗高、噪声大。
  \end{enumerate}
\end{enumerate}
\end{solution}

\begin{problem}
燃料燃烧一般分为哪几个阶段？各阶段的主要影响因素是什么？过量空气系数对锅炉燃烧的影响如何？
\end{problem}

\begin{solution}
\begin{enumerate}
  \item \textbf{燃烧阶段及主要影响因素}：
  \begin{enumerate}
    \item \textbf{预热阶段}：燃料预热、烘干和挥发分分解。燃料只从炉膛中吸取热量，并未开始燃烧，不需要供给空气。
    \item \textbf{燃烧阶段}：挥发分首先着火，随后焦炭剧烈燃烧放热。需要充足的空气供应（一次风满足挥发分燃烧，二次风满足焦炭燃烧）。
    \item \textbf{燃尽阶段}：焦炭质点被灰分包围，与空气接触不良，燃烧十分缓慢。所需空气量较少，但需要保证足够的时间让其燃尽。
  \end{enumerate}
  \item \textbf{过量空气系数的影响}：
  \begin{enumerate}
    \item \textbf{过量空气系数过大}：会因送入的冷空气过多而造成炉温降低，影响煤粉的着火燃烧，同时造成烟气容积增大，导致排烟热损失增大。
    \item \textbf{过量空气系数过小}：会因氧气不足而造成不完全燃烧热损失。
  \end{enumerate}
\end{enumerate}
\end{solution}

\begin{problem}
煤粉燃烧器一般有哪些类型？其结构原理及布置方式各有什么特点？
\end{problem}

\begin{solution}
\begin{enumerate}
  \item \textbf{类型}：按出口气流流态分为直流式和旋流式两大类。
  \item \textbf{直流燃烧器}：
  \begin{enumerate}
    \item \textbf{结构原理}：出口气流为直流射流，一、二次风以直流方式喷入炉膛。
    \item \textbf{布置方式特点}：布置在炉膛四角，形成切圆燃烧（如单切圆、双切圆等）。依靠切圆火炬的旋转和气流共同作用，一、二次风及热烟气混合强烈，形成有利的加热着火、燃烧及燃尽条件。
  \end{enumerate}
  \item \textbf{旋流燃烧器}：
  \begin{enumerate}
    \item \textbf{结构原理}：出口包含旋流射流，利用旋流叶片使气流旋转。
    \item \textbf{布置方式特点}：主要有前墙布置、两面墙对冲或交错布置、炉底或炉顶布置等多种方式。
  \end{enumerate}
\end{enumerate}
\end{solution}

\begin{problem}
何谓锅炉的水循环及自然水循环？自然循环锅炉与控制循环锅炉及直流锅炉的主要区别是什么？
\end{problem}

\begin{solution}
\begin{enumerate}
  \item \textbf{水循环及自然水循环}：锅炉水循环是指水和汽水混合物在蒸发设备循环回路中的连续流动。自然循环是利用下降管中的水和上升管中汽水混合物的密度差所产生的压头（水循环动力）来实现的流动。
  \item \textbf{主要区别}：
  \begin{enumerate}
    \item \textbf{自然循环锅炉}：有汽包，管内工质的流动完全依靠汽、水密度差产生的运动压头来实现。
    \item \textbf{控制循环锅炉}：有汽包，除了依靠汽水密度差外，主要在水循环回路中加装了锅水循环泵，利用泵的压头来克服流动阻力维持循环。
    \item \textbf{直流锅炉}：无汽包，受热面内工质均为强制流动且一次通过所有受热面，没有水循环，流动全靠给水泵所提供的压头。
  \end{enumerate}
\end{enumerate}
\end{solution}

\begin{problem}
锅炉汽水系统受热面包括哪几种？各受热面的作用如何？其结构及布置有什么不同？
\end{problem}

\begin{solution}
\begin{enumerate}
  \item \textbf{受热面种类}：省煤器、水冷壁、过热器、再热器。
  \item \textbf{作用、结构与布置}：
  \begin{enumerate}
    \item \textbf{省煤器（预热）}：利用尾部烟气预热给水，降低排烟温度，提高锅炉效率。一般为多排蛇形管，水平布置在尾部竖井中。
    \item \textbf{水冷壁（蒸发）}：吸收辐射热使水受热蒸发，是主要的蒸发受热面，并保护炉墙防结渣。由垂直管子组成，布置在炉膛四周内壁。
    \item \textbf{过热器（过热）}：将饱和蒸汽加热成具有一定过热度的过热蒸汽。结构为蛇形管或管屏；对流式布置在水平烟道或尾部竖井，半辐射式布置在炉膛出口上方，辐射式布置在炉膛内部或炉顶。
    \item \textbf{再热器（再热）}：将汽轮机高压缸排汽重新加热，管材需更耐高温。壁式辐射再热器为单排管布置在炉膛上部，其它多为管屏或蛇形管布置在对流过热器之后的水平烟道中。
  \end{enumerate}
\end{enumerate}
\end{solution}

\begin{problem}
空气预热器的作用有哪些？布置在锅炉的什么位置？在大容量锅炉中为什么一般采用回转式空气预热器？
\end{problem}

\begin{solution}
\begin{enumerate}
  \item \textbf{作用}：利用烟气的热量来加热燃烧所需的空气，它在烟气温度最低的区域回收了热量，降低了排烟温度，提高了锅炉效率，同时预热空气强化了着火燃烧过程。
  \item \textbf{布置位置}：布置在锅炉的最末端（即尾部烟道的低温区域）。
  \item \textbf{大容量锅炉采用回转式的原因}：回转式空气预热器（如三分仓结构）具有外形小、重量轻的优点，且其传热元件允许有较大的磨损。虽然它存在漏风大和结构复杂的缺点，但综合优势使其特别适用于大容量锅炉。
\end{enumerate}
\end{solution}
